\section{Integral means and the full integrability range}
\label{sec:integrability}

Proposition~\ref{prop:transfer-bound} and the disk comparison now
give the inverse-square estimate.  We first identify its sharp
integral-means exponent, then obtain the negative area exponents by
H\"older's inequality.  The positive exponents follow from a
separate weighted area estimate for univalent maps.

\subsection{The inverse-square endpoint and negative powers}

\begin{proof}[Proof of the integral-means assertion in Theorem~\ref{thm:main}]
 Proposition~\ref{prop:transfer-bound} gives $\rho\leq2$.
 Proposition~\ref{prop:disk-transfer} therefore implies that, for every
 $\epsilon>0$, there is $C_\epsilon<\infty$, independent of
 $f\in\SSS$, such that
 \begin{equation}\label{eq:final-inverse-means}
  M_{-2}[f'](r)\leq C_\epsilon(1-r)^{-1-\epsilon},
  \qquad \frac12\leq r<1.
 \end{equation}
 In particular $\beta_f(-2)\leq1$ for every $f\in\SSS$.

 For the reverse inequality consider the Koebe map
 \[
  f_{\mathrm K}(z)=\frac{z}{(1-z)^2},\qquad
  f_{\mathrm K}'(z)=\frac{1+z}{(1-z)^3}.
 \]
 It belongs to $\SSS$: equality of its values at $z,u\in\DD$
 gives $(z-u)(1-zu)=0$, hence $z=u$.
 Write $\delta=1-r$ and take $r$ sufficiently close to $1$.
 On the arc $z=re^{i(\pi+\theta)}$ with $|\theta|\leq\delta$,
 \[
  |1+z|\leq\delta+r|\theta|\leq2\delta,\qquad |1-z|\geq1.
 \]
 Consequently
 \[
  M_{-2}[f_{\mathrm K}'](r)
    \geq \frac1{2\pi}\int_{-\delta}^{\delta}
                      \frac{d\theta}{4\delta^2}
    =\frac1{4\pi\delta}.
 \]
 No exponent smaller than $1$ can bound these means near the
 boundary.  Hence $\beta_{f_{\mathrm K}}(-2)\geq1$, and
 $B_{\SSS}(-2)=1$.
\end{proof}

\begin{corollary}[The full negative linear tail]\label{cor:negative-linear-tail}
For every real $t\leq-2$ and every $\epsilon>0$, there is
$C_{t,\epsilon}<\infty$ such that
\begin{equation}\label{eq:negative-linear-tail}
 M_t[f'](r)\leq C_{t,\epsilon}(1-r)^{-(|t|-1)-\epsilon}
 \qquad\left(f\in\SSS,\ \frac12\leq r<1\right).
\end{equation}
Consequently,
\[
 B_{\SSS}(t)=|t|-1\qquad(t\leq-2).
\]
\end{corollary}

\begin{proof}
Put $a=|t|-2\geq0$.  The distortion lower bound
\eqref{prelim:disk-bounds} and \eqref{eq:final-inverse-means} give
\[
 \begin{aligned}
 M_t[f'](r)
 &\leq\left(\frac{(1+r)^3}{1-r}\right)^a M_{-2}[f'](r)\\
 &\leq 8^a C_\epsilon(1-r)^{-(a+1)-\epsilon}.
 \end{aligned}
\]
Since $a+1=|t|-1$, this proves \eqref{eq:negative-linear-tail} with
$C_{t,\epsilon}=8^{|t|-2}C_\epsilon$.  At $t=-2$ it is exactly
\eqref{eq:final-inverse-means}.  For each fixed $f$ and $t$, letting
$\epsilon\downarrow0$ in the definition of $\beta_f(t)$ gives
$\beta_f(t)\leq|t|-1$; taking the supremum over $f$ gives the same
upper bound for $B_{\SSS}(t)$.

For the Koebe map and the arc used above, $|f_{\mathrm K}'(z)|\leq2\delta$.
Thus, for $\delta=1-r$ sufficiently small,
\[
 M_t[f_{\mathrm K}'](r)
 \geq\frac1{2\pi}\int_{-\delta}^{\delta}(2\delta)^{-|t|}\,d\theta
 =\frac{2^{-|t|}}{\pi}\delta^{1-|t|}.
\]
Hence $\beta_{f_{\mathrm K}}(t)\geq|t|-1$, proving the equality.
\end{proof}

For $-2<t<0$, put $q=-t/2\in(0,1)$.  H\"older's inequality on the
probability circle gives
\begin{equation}\label{eq:negative-holder}
 M_t[f'](r)
   =\frac1{2\pi}\int_{-\pi}^{\pi}
        \bigl(|f'(re^{i\theta})|^{-2}\bigr)^q\dd\theta
   \leq M_{-2}[f'](r)^q.
\end{equation}
Choose $0<\epsilon<1/q-1$.  The exponent $q(1+\epsilon)$ in
\eqref{eq:final-inverse-means} is then less than $1$, so
\begin{equation}\label{eq:negative-area}
 \int_{\DD}|f'|^t\dd A
    =2\pi\int_0^1 r M_t[f'](r)\dd r<\infty
    \qquad(f\in\SSS,\ -2<t<0).
\end{equation}
There is no singularity on a compact subdisk because $f'$ never
vanishes.  The case $t=0$ gives the area $\pi$.

\subsection{Positive powers}

We use the classical bounds of Lemma~\ref{lem:classical-univalence}
to establish two elementary integrability statements for $f$ itself.

\begin{lemma}[Function means and area powers]\label{lem:function-powers}
 For every $0<p<1/2$ there is $C_p<\infty$ such that
 \begin{equation}\label{eq:function-means}
  \frac1{2\pi}\int_{-\pi}^{\pi}
       |f(re^{i\theta})|^p\dd\theta\leq C_p
  \qquad(f\in\SSS,\ 0<r<1).
 \end{equation}
 Moreover, for every $0\leq a<1$ and every $f\in\SSS$,
 \begin{equation}\label{eq:function-area}
  \int_{\DD}(1+|f(z)|)^a\dd A(z)<\infty.
 \end{equation}
\end{lemma}

\begin{proof}
 Away from the simple zero of $f$ at $0$,
 \[
  \Delta |f|^p=p^2|f|^{p-2}|f'|^2.
 \]
 The right side is locally integrable also at $0$, where it is
 $O(|z|^{p-2})$.  Green's radial mean identity gives
 \begin{equation}\label{eq:function-green}
 \frac1{2\pi}\int_{-\pi}^{\pi}|f(re^{i\theta})|^p\dd\theta
 =\frac{p^2}{2\pi}\int_{|z|<r}
   |f(z)|^{p-2}|f'(z)|^2\log\frac r{|z|}\dd A(z).
 \end{equation}
 To justify it at the zero, apply Green's theorem on an annulus
 \cite[Theorem~16.17]{Lee2013} with
 inner radius $\delta$.  The circular mean at that radius is
 $O(\delta^p)$ and its derivative is $O(\delta^{p-1})$, because
 $f(z)=z+O(z^2)$.  The inner terms are consequently
 $O(\delta^p(1+|\log\delta|))$ and vanish as $\delta\downarrow0$.

 Bound the logarithm in \eqref{eq:function-green} by
 $\log(1/|z|)$.  On $|z|\leq1/2$, the classical bounds give
 $|f(z)|\geq c|z|$ and $|f'(z)|\leq C$, uniformly in $\SSS$.
 Since $p-2<0$, this part is bounded by a constant times
 \[
  \int_0^{1/2}u^{p-1}\log(1/u)\dd u<\infty.
 \]
 On $|z|>1/2$, the growth bound
 $|f(z)|\leq |z|/(1-|z|)^2$ implies
 \[
  \log\frac1{|z|}
   \leq C(1-|z|)
   \leq C\min\{1,|f(z)|^{-1/2}\}.
 \]
 The conformal change of variables $w=f(z)$ therefore bounds the
 remaining integral by a constant times
 \[
  \int_{\CC}|w|^{p-2}\min\{1,|w|^{-1/2}\}\dd A(w)
  =2\pi\left(\int_0^1u^{p-1}\dd u
             +\int_1^\infty u^{p-3/2}\dd u\right)<\infty.
 \]
 This proves \eqref{eq:function-means}.

 For \eqref{eq:function-area}, the case $a=0$ is immediate.
 For $0<a<1$, choose
 \[
   \max\{0,a-\tfrac12\}<p<\tfrac12.
 \]
 If $a\leq p$, the power-mean inequality and
 \eqref{eq:function-means} bound the means of $|f|^a$.
 If $a>p$, the growth bound gives instead
 \[
  \frac1{2\pi}\int_{-\pi}^{\pi}|f(re^{i\theta})|^a\dd\theta
   \leq
   \left(\max_{|z|=r}|f(z)|\right)^{a-p}C_p
   \leq C_p(1-r)^{-2(a-p)}.
 \]
 Thus the means of $(1+|f|)^a$ are bounded by
 $C(1-r)^{-2(a-p)_+}$, where $(b)_+=\max\{b,0\}$.
 Since $2(a-p)_+<1$, radial integration proves the claim.
\end{proof}

\begin{lemma}[Positive derivative powers]\label{lem:positive-derivative-powers}
 For every $f\in\SSS$ and $0<t<2/3$,
 \[
  \int_{\DD}|f'(z)|^t\dd A(z)<\infty.
 \]
\end{lemma}

\begin{proof}
 Choose $0<\eta<2/t-3$ and set
 \[
  a=\frac{(2+\eta)t}{2-t}\in(0,1),\qquad
  u_\eta(z)=|f'(z)|^2(1+|f(z)|)^{-2-\eta}.
 \]
 Injective conformal area change gives
 \begin{equation}\label{eq:weighted-derivative-area}
  \int_{\DD}u_\eta\dd A
   =\int_{f(\DD)}(1+|w|)^{-2-\eta}\dd A(w)
   \leq\frac{2\pi}{\eta(1+\eta)}.
 \end{equation}
 Here the last constant is the polar integral over $\CC$.
 The exact factorization
 \[
  |f'|^t=u_\eta^{t/2}\bigl((1+|f|)^a\bigr)^{(2-t)/2}
 \]
 and H\"older's inequality with exponents $2/t$ and $2/(2-t)$
 now yield
 \[
  \int_{\DD}|f'|^t\dd A
   \leq
    \left(\int_{\DD}u_\eta\dd A\right)^{t/2}
    \left(\int_{\DD}(1+|f|)^a\dd A\right)^{(2-t)/2}<\infty
 \]
 by Lemma~\ref{lem:function-powers}.
\end{proof}

\subsection{Domains, equivalence, and endpoint sharpness}

\begin{proof}[Proof of the area assertion in Theorem~\ref{thm:main}]
 Equations~\eqref{eq:negative-area} and
 Lemma~\ref{lem:positive-derivative-powers}, together with $t=0$,
 prove the disk integral for $f\in\SSS$ throughout
 $-2<t<2/3$.  For an arbitrary univalent $f:\DD\to\CC$,
 \[
  g(z)=\frac{f(z)-f(0)}{f'(0)}\in\SSS,\qquad
  \int_{\DD}|f'|^t\dd A
    =|f'(0)|^t\int_{\DD}|g'|^t\dd A.
 \]
 The multiplier is finite and positive for every real $t$.

 If $\varphi:W\to\DD$ is a conformal bijection and
 $f=\varphi^{-1}$, the identities
 $\varphi'(f(w))=1/f'(w)$ and
 $\dd A(f(w))=|f'(w)|^2\dd A(w)$ give
 \begin{equation}\label{eq:domain-change}
  \int_W|\varphi'(z)|^s\dd A(z)
    =\int_{\DD}|f'(w)|^{2-s}\dd A(w).
 \end{equation}
 This change of variables is valid for nonnegative integrands,
 whether or not the integrals are finite.  Since
 $4/3<s<4$ is equivalent to $-2<2-s<2/3$, the desired domain
 integral is finite.
\end{proof}

For completeness, the disk and domain formulations have exactly
the same scope.  A univalent disk image is simply connected and
cannot equal $\CC$, since its inverse would be a bounded entire
nonconstant function.  If its boundary in the sphere had at most
one point, it would be a nonempty clopen subset of the connected
sphere with that possible point removed.  It would therefore be
the whole sphere or a punctured sphere, which is incompatible
with being a proper plane domain.  Thus it satisfies the stated
boundary-size condition, and \eqref{eq:domain-change} applies to
its inverse.  Conversely, each conformal bijection in the domain
formulation has a univalent inverse on $\DD$.

For a plane domain $U$ and $1\leq r\leq\infty$, write $L_r^1(U)$
for the homogeneous Sobolev space of locally integrable, weakly
differentiable real-valued functions with seminorm
$\|\nabla u\|_{L^r(U)}<\infty$, using the essential supremum when
$r=\infty$.

\begin{corollary}[Homogeneous Sobolev pullback]\label{cor:sobolev-pullback}
 Let $\Omega\subset\CC$ be a simply connected domain whose boundary
 in the Riemann sphere has at least two points, and let
 $\varphi:\Omega\to\DD$ be a
 conformal bijection. For $2<p<\infty$ and $4/3<s<4$, put
 $q=ps/(p+s-2)$. Then pullback $\varphi^*u=u\circ\varphi$ is bounded
 from $L_p^1(\DD)$ to $L_q^1(\Omega)$. It is also bounded from
 $L_\infty^1(\DD)$ to $L_s^1(\Omega)$ for every $4/3<s<4$.
 The finite bounds on gradient seminorms may depend on the map and
 the exponents.
\end{corollary}

\begin{proof}
 This is the known Sobolev formulation of Brennan's conjecture due to
 Gol'dshtein and Ukhlov
 \cite[Equivalence Theorem and its following Remark~2]{GoldshteinUkhlov2013}.
 Indeed, $1<q<p$ and $(p-2)q/(p-q)=s$, so Theorem~\ref{thm:main}
 supplies the derivative integrability required by their equivalence.
 The cited remark gives the case $p=\infty$, with $q=s$.
\end{proof}

The area assertion also implies the individual endpoint
integral-means bound.  Indeed, fix $f\in\SSS$ and $0<\eta<2$.
The function $|f'|^{-2+\eta}$ is subharmonic, since $f'$ is
nonvanishing; hence its circular means $M_{-2+\eta}[f'](r)$
are nondecreasing.  If its disk area integral is finite, then,
for $r\geq1/2$,
\[
 \int_{\DD}|f'|^{-2+\eta}\dd A
 \geq2\pi M_{-2+\eta}[f'](r)\int_r^1u\dd u,
\]
so $M_{-2+\eta}[f'](r)\leq C_{f,\eta}(1-r)^{-1}$.
The distortion lower bound gives
$|f'(re^{i\theta})|^{-\eta}\leq C_\eta(1-r)^{-\eta}$.
Multiplying these bounds yields
\[
 M_{-2}[f'](r)\leq C'_{f,\eta}(1-r)^{-1-\eta}.
\]
Letting $\eta\downarrow0$ in the definition of the exponent gives
$\beta_f(-2)\leq1$.  In the other direction, this individual
exponent bound gives an estimate with every positive exponent
loss, and \eqref{eq:negative-holder} gives all negative interior
area exponents.  The positive exponents were proved independently
above.  Together with the Koebe lower bound, this establishes the
stated equivalence with $B_{\SSS}(-2)=1$.  In our proof, the uniform constant in
\eqref{eq:final-inverse-means} comes directly from the transfer-operator
argument.

\begin{proof}[Proof of Corollary~\ref{cor:sharp-endpoints}]
 Use the same Koebe map $f_{\mathrm K}$.  For sufficiently small
 $|w|$,
 \[
  |f_{\mathrm K}'(-1+w)|=\frac{|w|}{|2-w|^3}\asymp |w|,
  \qquad
  |f_{\mathrm K}'(1-w)|=\frac{|2-w|}{|w|^3}\asymp |w|^{-3}.
 \]
 The fixed sector
 \[
  0<|w|<\tfrac12,\qquad |\arg w|<\tfrac{\pi}{3}
 \]
 lies inside $\DD$ under either substitution $z=-1+w$ or
 $z=1-w$: in both cases $|z|^2=1-2\Re w+|w|^2<1$.
 The area integral at $t=-2$ near $-1$ and the one at $t=2/3$
 near $1$ therefore each dominate a positive multiple of
 \[
  \int_0^{1/2}u^{-2}u\dd u=\infty.
 \]
 More generally these sectors diverge for $t\leq-2$ and
 $t\geq2/3$, respectively.  The image of $f_{\mathrm K}$ meets
 the domain hypotheses just verified, so its inverse, by
 \eqref{eq:domain-change}, also gives divergence at $s=4$ and
 $s=4/3$.  Thus both open endpoints are sharp.
\end{proof}
